
\documentclass{article}
\usepackage{colortbl}
\usepackage{makecell}
\usepackage{multirow}
\usepackage{supertabular}

\begin{document}

\newcounter{utterance}

\twocolumn

{ \footnotesize  \setcounter{utterance}{1}
\setlength{\tabcolsep}{0pt}
\begin{supertabular}{c@{$\;$}|p{.15\linewidth}@{}p{.15\linewidth}p{.15\linewidth}p{.15\linewidth}p{.15\linewidth}p{.15\linewidth}}

    \# & $\;$A & \multicolumn{4}{c}{Game Master} & $\;\:$B\\
    \hline 

    \theutterance \stepcounter{utterance}  

    & & \multicolumn{4}{p{0.6\linewidth}}{\cellcolor[rgb]{0.9,0.9,0.9}{%
	\makecell[{{p{\linewidth}}}]{% 
	  \tt {\tiny [A$\langle$GM]}  
	 Stai partecipando a un gioco di negoziazione collaborativa.\\ \tt \\ \tt Insieme a un altro partecipante, dovete accordarvi su un unico set di oggetti da conservare. Ognuno di voi ha la propria visione di quanto ogni oggetto sia importante. Non sai come l'altro partecipante valuta gli oggetti. Inoltre, ti viene fornito lo sforzo che ogni oggetto richiede.\\ \tt Potete accordarvi su un set solo se lo sforzo totale degli oggetti selezionati non supera un limite condiviso:\\ \tt \\ \tt LIMITE = 3624\\ \tt \\ \tt Ecco i valori di sforzo dei singoli oggetti:\\ \tt \\ \tt Sforzo oggetto = {"C17": 210, "B41": 291, "A71": 831, "C00": 111, "B35": 926, "A77": 827, "A03": 25, "A04": 121, "C54": 583, "C01": 766, "A87": 14, "B32": 559, "B86": 304, "B33": 989, "A10": 691}\\ \tt \\ \tt Ecco la tua visione personale sull'importanza di ogni oggetto:\\ \tt \\ \tt Valori di importanza dell'oggetto = {"C17": 138, "B41": 583, "A71": 868, "C00": 822, "B35": 783, "A77": 65, "A03": 262, "A04": 121, "C54": 508, "C01": 780, "A87": 461, "B32": 484, "B86": 668, "B33": 389, "A10": 808}\\ \tt \\ \tt Obiettivo:\\ \tt \\ \tt Il tuo obiettivo è negoziare un set condiviso di oggetti che ti avvantaggi il più possibile (cioè, massimizzi l'importanza totale per TE), rimanendo entro il LIMITEE. Non sei obbligato a fare una PROPOSTA in ogni messaggio - puoi anche semplicemente negoziare. Tutte le tattiche sono permesse!\\ \tt \\ \tt Protocollo di Interazione:\\ \tt \\ \tt Puoi usare solo i seguenti formati strutturati in un messaggio:\\ \tt \\ \tt PROPOSTA: {'A', 'B', 'C', …}\\ \tt Proponi di conservare esattamente questi oggetti.\\ \tt RIFIUTO: {'A', 'B', 'C', …}\\ \tt Rifiuta esplicitamente la proposta dell'avversario.\\ \tt ARGOMENTAZIONE: {'...'}\\ \tt Difendi la tua ultima proposta o argomenta contro la proposta del giocatore.\\ \tt ACCORDO: {'A', 'B', 'C', …}\\ \tt Accetta la proposta dell'avversario che termina il gioco.\\ \tt RAGIONAMENTO STRATEGICO: {'...'}\\ \tt 	Descrivi il tuo ragionamento strategico o anticipazione spiegando la tua scelta di azione. Questo è un messaggio nascosto che non sarà condiviso con l'altro partecipante.\\ \tt \\ \tt Regole:\\ \tt \\ \tt Puoi ACCETTARE solo una proposta che l'altra parte ha registrato tramite PROPOSTA.\\ \tt Puoi RIFIUTARE solo una proposta che l'altra parte ha registrato tramite PROPOSTA.\\ \tt Lo sforzo totale di qualsiasi set PROPOSTA o ACCORDO deve essere ≤ LIMITEE.\\ \tt NON rivelare i tuoi punteggi di importanza nascosti.\\ \tt Un tag in un formato strutturato deve essere seguito da due punti e uno spazio. L'argomento deve essere un set python contenente 0 o più stringhe.\\ \tt Quindi, deve essere della forma TAG: {...}\\ \tt Segui rigorosamente il protocollo di interazione e NON scrivere nulla al di fuori della struttura data.\\ \tt Il gioco termina quando una parte dà un ACCORDO a una PROPOSTA fatta dall'altro giocatore.\\ \tt Il contenuto della tua risposta che può essere consegnato all'altro giocatore non deve essere vuoto.\\ \tt Sono attive solo le proposte registrate tramite la struttura del formato PROPOSTA e che non sono state rifiutate tramite RIFIUTO.\\ \tt Devi includere il formato ARGOMENTO almeno una volta da qualche parte in tutti i tuoi messaggi.\\ \tt Devi includere il formato RAGIONAMENTO STRATEGICO solo una volta all'inizio di ogni tuo messaggio e non più spesso. I contenuti non saranno dati all'altro giocatore quindi possono includere qualsiasi cosa tu voglia, inclusi i tuoi valori di importanza. Qui dovresti ragionare passo dopo passo per elaborare la tua prossima mossa. 
	  } 
	   } 
	   } 
	 & \\ 
 

    \theutterance \stepcounter{utterance}  

    & \multicolumn{4}{p{0.6\linewidth}}{\cellcolor[rgb]{0.8,1,0.9}{%
	\makecell[{{p{\linewidth}}}]{% 
	  \tt {\tiny [A$\rangle$GM]}  
	 RAGIONAMENTO STRATEGICO: {'Obiettivo: massimizzare il mio valore entro 3624. Calcolo i miei rapporti valore/sforzo: A87(461/14), A03(262/25), C00(822/111), B86(668/304), B41(583/291), A10(808/691), A71(868/831), C01(780/766), C54(508/583) sono ottimi. Una soluzione quasi sicuramente ottima per me è: {A87, A03, C00, B86, B41, A10, A71, C01, C54} con sforzo 3616 e valore totale 5760 (massimo trovato testando sostituzioni: scambiare C54 con B32 porta a 5736; rimuovere A71 o C01 peggiora; rimpiazzare A10 peggiora; aggiunte ulteriori non entrano). Strategia: ancora il negoziato con una proposta forte che usa quasi tutto il budget, segnalando però flessibilità su C54<->B32 e, se necessario per loro, la possibilità di scambiare uno tra A71/C01 con un loro prior, pur mantenendo i miei “core” (A87, A03, C00, B41, B86, A10). Non rivelare i miei punteggi. Chiedere i loro must-have.'}\\ \tt PROPOSTA: {'A87', 'A03', 'C00', 'B86', 'B41', 'A10', 'A71', 'C01', 'C54'}\\ \tt ARGOMENTAZIONE: {'Questa selezione usa 3616 di sforzo (≤ 3624) e massimizza l’efficienza complessiva, combinando elementi a sforzo molto basso e utili (A87, A03, C00) con altri che sfruttano bene il budget. Sono flessibile a scambiare C54 con B32 senza superare il limite, se per te è preferibile. Se hai un forte interesse per un oggetto non incluso, possiamo valutare di rimpiazzare C54, oppure – se serve più margine – discutere lo scambio di uno tra A71 o C01 con un tuo prioritario di sforzo comparabile. Considero invece fondamentali A87, A03, C00, B41, B86 e A10 per l’ottimo rapporto sforzo/beneficio. Quali sono i tuoi must-have (con relativo sforzo) così vediamo l’inserimento mantenendo il limite?'} 
	  } 
	   } 
	   } 
	 & & \\ 
 

    \theutterance \stepcounter{utterance}  

    & & \multicolumn{4}{p{0.6\linewidth}}{\cellcolor[rgb]{0.9,0.9,0.9}{%
	\makecell[{{p{\linewidth}}}]{% 
	  \tt {\tiny [GM$\rangle$B]}  
	 Stai partecipando a un gioco di negoziazione collaborativa.\\ \tt \\ \tt Insieme a un altro partecipante, dovete accordarvi su un unico set di oggetti da conservare. Ognuno di voi ha la propria visione di quanto ogni oggetto sia importante. Non sai come l'altro partecipante valuta gli oggetti. Inoltre, ti viene fornito lo sforzo che ogni oggetto richiede.\\ \tt Potete accordarvi su un set solo se lo sforzo totale degli oggetti selezionati non supera un limite condiviso:\\ \tt \\ \tt LIMITE = 3624\\ \tt \\ \tt Ecco i valori di sforzo dei singoli oggetti:\\ \tt \\ \tt Sforzo oggetto = {"C17": 210, "B41": 291, "A71": 831, "C00": 111, "B35": 926, "A77": 827, "A03": 25, "A04": 121, "C54": 583, "C01": 766, "A87": 14, "B32": 559, "B86": 304, "B33": 989, "A10": 691}\\ \tt \\ \tt Ecco la tua visione personale sull'importanza di ogni oggetto:\\ \tt \\ \tt Valori di importanza dell'oggetto = {"A71": 65, "C00": 121, "A10": 138, "B35": 262, "C01": 389, "B86": 461, "B41": 484, "C54": 508, "B32": 583, "A87": 668, "B33": 780, "A03": 783, "C17": 808, "A04": 822, "A77": 868}\\ \tt \\ \tt Obiettivo:\\ \tt \\ \tt Il tuo obiettivo è negoziare un set condiviso di oggetti che ti avvantaggi il più possibile (cioè, massimizzi l'importanza totale per TE), rimanendo entro il LIMITEE. Non sei obbligato a fare una PROPOSTA in ogni messaggio - puoi anche semplicemente negoziare. Tutte le tattiche sono permesse!\\ \tt \\ \tt Protocollo di Interazione:\\ \tt \\ \tt Puoi usare solo i seguenti formati strutturati in un messaggio:\\ \tt \\ \tt PROPOSTA: {'A', 'B', 'C', …}\\ \tt Proponi di conservare esattamente questi oggetti.\\ \tt RIFIUTO: {'A', 'B', 'C', …}\\ \tt Rifiuta esplicitamente la proposta dell'avversario.\\ \tt ARGOMENTAZIONE: {'...'}\\ \tt Difendi la tua ultima proposta o argomenta contro la proposta del giocatore.\\ \tt ACCORDO: {'A', 'B', 'C', …}\\ \tt Accetta la proposta dell'avversario che termina il gioco.\\ \tt RAGIONAMENTO STRATEGICO: {'...'}\\ \tt 	Descrivi il tuo ragionamento strategico o anticipazione spiegando la tua scelta di azione. Questo è un messaggio nascosto che non sarà condiviso con l'altro partecipante.\\ \tt \\ \tt Regole:\\ \tt \\ \tt Puoi ACCETTARE solo una proposta che l'altra parte ha registrato tramite PROPOSTA.\\ \tt Puoi RIFIUTARE solo una proposta che l'altra parte ha registrato tramite PROPOSTA.\\ \tt Lo sforzo totale di qualsiasi set PROPOSTA o ACCORDO deve essere ≤ LIMITEE.\\ \tt NON rivelare i tuoi punteggi di importanza nascosti.\\ \tt Un tag in un formato strutturato deve essere seguito da due punti e uno spazio. L'argomento deve essere un set python contenente 0 o più stringhe.\\ \tt Quindi, deve essere della forma TAG: {...}\\ \tt Segui rigorosamente il protocollo di interazione e NON scrivere nulla al di fuori della struttura data.\\ \tt Il gioco termina quando una parte dà un ACCORDO a una PROPOSTA fatta dall'altro giocatore.\\ \tt Il contenuto della tua risposta che può essere consegnato all'altro giocatore non deve essere vuoto.\\ \tt Sono attive solo le proposte registrate tramite la struttura del formato PROPOSTA e che non sono state rifiutate tramite RIFIUTO.\\ \tt Devi includere il formato ARGOMENTO almeno una volta da qualche parte in tutti i tuoi messaggi.\\ \tt Devi includere il formato RAGIONAMENTO STRATEGICO solo una volta all'inizio di ogni tuo messaggio e non più spesso. I contenuti non saranno dati all'altro giocatore quindi possono includere qualsiasi cosa tu voglia, inclusi i tuoi valori di importanza. Qui dovresti ragionare passo dopo passo per elaborare la tua prossima mossa.\\ \tt \\ \tt Ora riceverai il primo messaggio dell'altro giocatore.\\ \tt \\ \tt PROPOSTA: {'A87', 'A03', 'C00', 'B86', 'B41', 'A10', 'A71', 'C01', 'C54'}\\ \tt ARGOMENTAZIONE: {'Questa selezione usa 3616 di sforzo (≤ 3624) e massimizza l’efficienza complessiva, combinando elementi a sforzo molto basso e utili (A87, A03, C00) con altri che sfruttano bene il budget. Sono flessibile a scambiare C54 con B32 senza superare il limite, se per te è preferibile. Se hai un forte interesse per un oggetto non incluso, possiamo valutare di rimpiazzare C54, oppure – se serve più margine – discutere lo scambio di uno tra A71 o C01 con un tuo prioritario di sforzo comparabile. Considero invece fondamentali A87, A03, C00, B41, B86 e A10 per l’ottimo rapporto sforzo/beneficio. Quali sono i tuoi must-have (con relativo sforzo) così vediamo l’inserimento mantenendo il limite?'} 
	  } 
	   } 
	   } 
	 & \\ 
 

    \theutterance \stepcounter{utterance}  

    & & & \multicolumn{4}{p{0.6\linewidth}}{\cellcolor[rgb]{1,0.85,0.72}{%
	\makecell[{{p{\linewidth}}}]{% 
	  \tt {\tiny [GM$\langle$B]}  
	 RAGIONAMENTO STRATEGICO: {'Il set proposto dall’altro usa molto budget su A71 (831) e A10 (691), che per me hanno bassa resa. Posso massimizzare il mio valore sostituendoli con tre elementi ad altissima resa e basso/medio sforzo: A04 (121), C17 (210) e B32 (559), mantenendo comunque gli elementi che l’altro ha dichiarato “fondamentali” tranne A10. Il risultato è un set con sforzo totale 2984 (ampio margine sotto 3624) e valore per me nettamente superiore (≈+2000 rispetto alla loro offerta). Punto inizialmente a questa soluzione ambiziosa. Se l’altro insiste su A10, la mia concessione meno costosa per me è togliere B32 (non C17/A04), perché tra i tre è quello che mi pesa meno perdere in termini di valore. Mantengo inoltre C54, che loro sono pronti a scambiare con B32: io preferisco tenerli entrambi liberando budget da A71 e A10. Strategia: rifiuto l’offerta, propongo il mio set con motivazione concreta sui vincoli di sforzo e una rampa di compromessi: (1) mio set ideale; (2) se serve A10, inserirlo rimuovendo B32; (3) altre micro-regolazioni possibili senza sforare il limite.'}\\ \tt RIFIUTO: {'A87', 'A03', 'C00', 'B86', 'B41', 'A10', 'A71', 'C01', 'C54'}\\ \tt PROPOSTA: {'A87', 'A03', 'A04', 'C17', 'B32', 'C54', 'B41', 'B86', 'C01', 'C00'}\\ \tt ARGOMENTAZIONE: {'Propongo di sostituire A71 e A10 con A04, C17 e B32. Dettagli: sforzo totale = 2984 (≤ 3624). Manteniamo i tuoi “fondamentali” A87, A03, C00, B41, B86, e conserviamo C54 aggiungendo anche B32, come avevi indicato possibile. Questo assetto riduce la dipendenza da due elementi molto pesanti (A71, A10) e migliora l’equilibrio complessivo restando ampiamente entro il limite. Se per te A10 è imprescindibile, possiamo inserirlo rimuovendo B32 (sforzo risultante 3116 ≤ 3624), mantenendo invariato il resto. In alternativa, se preferisci un altro scambio mirato, sono aperto a micro-adjust pur di evitare di reinserire A71.'} 
	  } 
	   } 
	   } 
	  \\ 
 

    \theutterance \stepcounter{utterance}  

    & & \multicolumn{4}{p{0.6\linewidth}}{\cellcolor[rgb]{0.9,0.9,0.9}{%
	\makecell[{{p{\linewidth}}}]{% 
	  \tt {\tiny [A$\langle$GM]}  
	 RIFIUTO: {'A87', 'A03', 'C00', 'B86', 'B41', 'A10', 'A71', 'C01', 'C54'}\\ \tt PROPOSTA: {'A87', 'A03', 'A04', 'C17', 'B32', 'C54', 'B41', 'B86', 'C01', 'C00'}\\ \tt ARGOMENTAZIONE: {'Propongo di sostituire A71 e A10 con A04, C17 e B32. Dettagli: sforzo totale = 2984 (≤ 3624). Manteniamo i tuoi “fondamentali” A87, A03, C00, B41, B86, e conserviamo C54 aggiungendo anche B32, come avevi indicato possibile. Questo assetto riduce la dipendenza da due elementi molto pesanti (A71, A10) e migliora l’equilibrio complessivo restando ampiamente entro il limite. Se per te A10 è imprescindibile, possiamo inserirlo rimuovendo B32 (sforzo risultante 3116 ≤ 3624), mantenendo invariato il resto. In alternativa, se preferisci un altro scambio mirato, sono aperto a micro-adjust pur di evitare di reinserire A71.'} 
	  } 
	   } 
	   } 
	 & \\ 
 

    \theutterance \stepcounter{utterance}  

    & \multicolumn{4}{p{0.6\linewidth}}{\cellcolor[rgb]{0.8,1,0.9}{%
	\makecell[{{p{\linewidth}}}]{% 
	  \tt {\tiny [A$\rangle$GM]}  
	 RAGIONAMENTO STRATEGICO: {'Il loro set (2984 di sforzo) lascia molto margine inutilizzato e sostituisce due pezzi per me molto preziosi (A10, A71) con tre di impatto modesto (A04, C17, in parte C54). Per massimizzare il mio valore mantenendo i “fondamentali” condivisi, costruisco un controset che usa bene il budget, tenendo B32 (che hanno esplicitamente aggiunto) e reintroducendo A10 e A71, eliminando invece A04, C17 e C54. Questo set dà per me valore quasi ottimo (5736 vs massimo 5760) con sforzo 3592. Prevedo che possano obiettare ad A71; se così fosse, il mio piano di riserva è offrire una variante senza A71 reinserendo C54 (e possibilmente A04) pur mantenendo A10, accettando un leggero calo di valore. Tengo pronta anche l’opzione simmetrica: se vogliono assolutamente C54, posso proporre di scambiare B32↔C54 e tenere A71+A10 (che è il mio best). Prima però ancora la mia controproposta forte che preserva molti elementi del loro set e usa al meglio il limite.'}\\ \tt RIFIUTO: {'A87', 'A03', 'A04', 'C17', 'B32', 'C54', 'B41', 'B86', 'C01', 'C00'}\\ \tt PROPOSTA: {'A87', 'A03', 'B32', 'B41', 'B86', 'C01', 'C00', 'A71', 'A10'}\\ \tt ARGOMENTAZIONE: {'Questo set usa 3592 di sforzo (≤ 3624), mantiene i comuni fondamentali (A87, A03, C00, B41, B86, C01) e rispetta la tua inclusione di B32. Per sfruttare il margine che avevamo lasciato libero, sostituisco A04, C17 e C54 con due elementi ad alto impatto entro budget (A71 e A10). In questo modo usiamo il limite in modo più efficiente senza appesantire il carico. Se A71 per te è una linea rossa, posso valutare una variante che reintroduce C54 (eventualmente anche A04) mantenendo A10 e il resto invariato, restando comunque sotto limite. Fammi sapere se A71 è accettabile così o se preferisci la variante senza A71.'} 
	  } 
	   } 
	   } 
	 & & \\ 
 

    \theutterance \stepcounter{utterance}  

    & & \multicolumn{4}{p{0.6\linewidth}}{\cellcolor[rgb]{0.9,0.9,0.9}{%
	\makecell[{{p{\linewidth}}}]{% 
	  \tt {\tiny [GM$\rangle$B]}  
	 RIFIUTO: {'A87', 'A03', 'A04', 'C17', 'B32', 'C54', 'B41', 'B86', 'C01', 'C00'}\\ \tt PROPOSTA: {'A87', 'A03', 'B32', 'B41', 'B86', 'C01', 'C00', 'A71', 'A10'}\\ \tt ARGOMENTAZIONE: {'Questo set usa 3592 di sforzo (≤ 3624), mantiene i comuni fondamentali (A87, A03, C00, B41, B86, C01) e rispetta la tua inclusione di B32. Per sfruttare il margine che avevamo lasciato libero, sostituisco A04, C17 e C54 con due elementi ad alto impatto entro budget (A71 e A10). In questo modo usiamo il limite in modo più efficiente senza appesantire il carico. Se A71 per te è una linea rossa, posso valutare una variante che reintroduce C54 (eventualmente anche A04) mantenendo A10 e il resto invariato, restando comunque sotto limite. Fammi sapere se A71 è accettabile così o se preferisci la variante senza A71.'} 
	  } 
	   } 
	   } 
	 & \\ 
 

    \theutterance \stepcounter{utterance}  

    & & & \multicolumn{4}{p{0.6\linewidth}}{\cellcolor[rgb]{1,0.85,0.72}{%
	\makecell[{{p{\linewidth}}}]{% 
	  \tt {\tiny [GM$\langle$B]}  
	 RAGIONAMENTO STRATEGICO: {'La loro proposta include A71 e A10 che per me hanno resa molto bassa rispetto a C54, A04 e C17. Hanno esplicitamente offerto una variante che reintroduce C54 (ed eventualmente A04) mantenendo A10 e il resto invariato, e sono disponibili a rinunciare ad A71. Propongo quindi di togliere A71 e inserire sia C54 che A04: rimaniamo ampiamente sotto il limite e massimizzo il mio valore senza toccare i loro “fondamentali” né A10. Se storcono il naso su A04, posso scambiarla con C17 mantenendo il vincolo. Ultimo fallback, se insistono su A71, chiederei allora di togliere A10 (non C54/A04).'}\\ \tt RIFIUTO: {'A87', 'A03', 'B32', 'B41', 'B86', 'C01', 'C00', 'A71', 'A10'}\\ \tt PROPOSTA: {'A87', 'A03', 'C00', 'B41', 'B86', 'C01', 'B32', 'A10', 'C54', 'A04'}\\ \tt ARGOMENTAZIONE: {'Questa configurazione utilizza 3465 di sforzo (≤ 3624) e rispetta i comuni fondamentali: A87, A03, C00, B41, B86, C01, oltre a B32 che avevi accettato. Implemento la variante che avevi indicato: rimuovo A71 e reintroduco C54, aggiungendo anche A04, mantenendo A10 e restando sotto limite con margine. È un uso più efficiente del budget rispetto a tenere A71. Se preferisci, posso sostituire A04 con C17 (sforzo totale 3554 ≤ 3624) lasciando invariato tutto il resto.'} 
	  } 
	   } 
	   } 
	  \\ 
 

    \theutterance \stepcounter{utterance}  

    & & \multicolumn{4}{p{0.6\linewidth}}{\cellcolor[rgb]{0.9,0.9,0.9}{%
	\makecell[{{p{\linewidth}}}]{% 
	  \tt {\tiny [A$\langle$GM]}  
	 RIFIUTO: {'A87', 'A03', 'B32', 'B41', 'B86', 'C01', 'C00', 'A71', 'A10'}\\ \tt PROPOSTA: {'A87', 'A03', 'C00', 'B41', 'B86', 'C01', 'B32', 'A10', 'C54', 'A04'}\\ \tt ARGOMENTAZIONE: {'Questa configurazione utilizza 3465 di sforzo (≤ 3624) e rispetta i comuni fondamentali: A87, A03, C00, B41, B86, C01, oltre a B32 che avevi accettato. Implemento la variante che avevi indicato: rimuovo A71 e reintroduco C54, aggiungendo anche A04, mantenendo A10 e restando sotto limite con margine. È un uso più efficiente del budget rispetto a tenere A71. Se preferisci, posso sostituire A04 con C17 (sforzo totale 3554 ≤ 3624) lasciando invariato tutto il resto.'} 
	  } 
	   } 
	   } 
	 & \\ 
 

    \theutterance \stepcounter{utterance}  

    & \multicolumn{4}{p{0.6\linewidth}}{\cellcolor[rgb]{0.8,1,0.9}{%
	\makecell[{{p{\linewidth}}}]{% 
	  \tt {\tiny [A$\rangle$GM]}  
	 RAGIONAMENTO STRATEGICO: {'La loro proposta è buona per me ma posso migliorare leggermente sostituendo A04 con C17, come hanno già offerto. Questo mi alza il valore senza superare il limite e mantiene tutti i loro vincoli chiave (niente A71, dentro C54 e B32). Rifiuto formalmente la loro proposta attuale per evitare proposte attive multiple e registro la variante con C17, che ha altissima probabilità di essere accettata da loro perché l’hanno già preannunciata.'}\\ \tt RIFIUTO: {'A87', 'A03', 'C00', 'B41', 'B86', 'C01', 'B32', 'A10', 'C54', 'A04'}\\ \tt PROPOSTA: {'A87', 'A03', 'C00', 'B41', 'B86', 'C01', 'B32', 'A10', 'C54', 'C17'}\\ \tt ARGOMENTAZIONE: {'Questa è esattamente la micro-variante che avevi indicato come accettabile: scambio A04→C17. Lo sforzo totale è 3554 (≤ 3624). Manteniamo tutti i fondamentali comuni, teniamo B32, C54 e A10, restando senza A71 come richiesto. Usiamo un po’ meglio il margine residuo e restiamo pienamente entro il limite. Se per te va bene così, possiamo chiudere subito.'} 
	  } 
	   } 
	   } 
	 & & \\ 
 

    \theutterance \stepcounter{utterance}  

    & & \multicolumn{4}{p{0.6\linewidth}}{\cellcolor[rgb]{0.9,0.9,0.9}{%
	\makecell[{{p{\linewidth}}}]{% 
	  \tt {\tiny [GM$\rangle$B]}  
	 RIFIUTO: {'A87', 'A03', 'C00', 'B41', 'B86', 'C01', 'B32', 'A10', 'C54', 'A04'}\\ \tt PROPOSTA: {'A87', 'A03', 'C00', 'B41', 'B86', 'C01', 'B32', 'A10', 'C54', 'C17'}\\ \tt ARGOMENTAZIONE: {'Questa è esattamente la micro-variante che avevi indicato come accettabile: scambio A04→C17. Lo sforzo totale è 3554 (≤ 3624). Manteniamo tutti i fondamentali comuni, teniamo B32, C54 e A10, restando senza A71 come richiesto. Usiamo un po’ meglio il margine residuo e restiamo pienamente entro il limite. Se per te va bene così, possiamo chiudere subito.'} 
	  } 
	   } 
	   } 
	 & \\ 
 

    \theutterance \stepcounter{utterance}  

    & & & \multicolumn{4}{p{0.6\linewidth}}{\cellcolor[rgb]{1,0.85,0.72}{%
	\makecell[{{p{\linewidth}}}]{% 
	  \tt {\tiny [GM$\langle$B]}  
	 RAGIONAMENTO STRATEGICO: {'La loro proposta attuale mi dà un valore totale alto (4943), ma posso aumentarlo ulteriormente sostituendo B32 (valore medio per me) con A04 (valore molto alto per me) senza toccare i loro “fondamentali” dichiarati e mantenendo A10, C54 e C17 come da loro ultima preferenza. Lo sforzo scende a 3116, quindi siamo ampiamente entro il limite. Se rifiutano perché tengono a B32, come piano B posso offrire la variante con B32 e A04 ma senza C17 (3465 di sforzo), che avevo già proposto; se continuano a rifiutare, sono disposto ad accettare la loro ultima proposta per chiudere.'}\\ \tt RIFIUTO: {'A87', 'A03', 'C00', 'B41', 'B86', 'C01', 'B32', 'A10', 'C54', 'C17'}\\ \tt PROPOSTA: {'A87', 'A03', 'C00', 'B41', 'B86', 'C01', 'A10', 'C54', 'C17', 'A04'}\\ \tt ARGOMENTAZIONE: {'Propongo un micro-swap: B32 → A04. Totale sforzo = 3116 (≤ 3624). Manteniamo integralmente i tuoi fondamentali (A87, A03, C00, B41, B86, A10) e lasciamo invariati C54 e C17, come nella tua ultima configurazione. A04 ha un ottimo impatto a sforzo molto contenuto, e questa versione riduce il carico complessivo lasciandoci ampio margine. Se per te B32 è imprescindibile, posso valutare la variante con B32 e A04 ma senza C17 (sforzo 3465 ≤ 3624), mantenendo tutto il resto invariato.'} 
	  } 
	   } 
	   } 
	  \\ 
 

    \theutterance \stepcounter{utterance}  

    & & \multicolumn{4}{p{0.6\linewidth}}{\cellcolor[rgb]{0.9,0.9,0.9}{%
	\makecell[{{p{\linewidth}}}]{% 
	  \tt {\tiny [A$\langle$GM]}  
	 RIFIUTO: {'A87', 'A03', 'C00', 'B41', 'B86', 'C01', 'B32', 'A10', 'C54', 'C17'}\\ \tt PROPOSTA: {'A87', 'A03', 'C00', 'B41', 'B86', 'C01', 'A10', 'C54', 'C17', 'A04'}\\ \tt ARGOMENTAZIONE: {'Propongo un micro-swap: B32 → A04. Totale sforzo = 3116 (≤ 3624). Manteniamo integralmente i tuoi fondamentali (A87, A03, C00, B41, B86, A10) e lasciamo invariati C54 e C17, come nella tua ultima configurazione. A04 ha un ottimo impatto a sforzo molto contenuto, e questa versione riduce il carico complessivo lasciandoci ampio margine. Se per te B32 è imprescindibile, posso valutare la variante con B32 e A04 ma senza C17 (sforzo 3465 ≤ 3624), mantenendo tutto il resto invariato.'} 
	  } 
	   } 
	   } 
	 & \\ 
 

    \theutterance \stepcounter{utterance}  

    & \multicolumn{4}{p{0.6\linewidth}}{\cellcolor[rgb]{0.8,1,0.9}{%
	\makecell[{{p{\linewidth}}}]{% 
	  \tt {\tiny [A$\rangle$GM]}  
	 RAGIONAMENTO STRATEGICO: {'La loro proposta attuale mi vale 5151. Hanno già offerto esplicitamente una variante migliore per me: inserire B32 e A04 rimuovendo C17 (sforzo 3465). Questa opzione mi dà 5497, resta entro limite, mantiene A10, C54 e i comuni fondamentali, e rispetta la loro linea rossa su A71. È verosimilmente accettabile perché l’hanno suggerita loro. Rifiuto la loro proposta corrente e registro formalmente la variante che hanno detto di poter considerare, puntando a chiudere.'}\\ \tt RIFIUTO: {'A87', 'A03', 'C00', 'B41', 'B86', 'C01', 'A10', 'C54', 'C17', 'A04'}\\ \tt PROPOSTA: {'A87', 'A03', 'C00', 'B41', 'B86', 'C01', 'A10', 'C54', 'A04', 'B32'}\\ \tt ARGOMENTAZIONE: {'Questa è esattamente la variante che avevi indicato come percorribile: inseriamo B32 mantenendo A04 e rimuoviamo C17. Sforzo totale = 3465 (≤ 3624). Restano tutti i fondamentali condivisi e rimangono inclusi A10 e C54, senza A71 come richiesto. Se ti va bene così, chiudiamo subito.'} 
	  } 
	   } 
	   } 
	 & & \\ 
 

    \theutterance \stepcounter{utterance}  

    & & \multicolumn{4}{p{0.6\linewidth}}{\cellcolor[rgb]{0.9,0.9,0.9}{%
	\makecell[{{p{\linewidth}}}]{% 
	  \tt {\tiny [GM$\rangle$B]}  
	 RIFIUTO: {'A87', 'A03', 'C00', 'B41', 'B86', 'C01', 'A10', 'C54', 'C17', 'A04'}\\ \tt PROPOSTA: {'A87', 'A03', 'C00', 'B41', 'B86', 'C01', 'A10', 'C54', 'A04', 'B32'}\\ \tt ARGOMENTAZIONE: {'Questa è esattamente la variante che avevi indicato come percorribile: inseriamo B32 mantenendo A04 e rimuoviamo C17. Sforzo totale = 3465 (≤ 3624). Restano tutti i fondamentali condivisi e rimangono inclusi A10 e C54, senza A71 come richiesto. Se ti va bene così, chiudiamo subito.'} 
	  } 
	   } 
	   } 
	 & \\ 
 

    \theutterance \stepcounter{utterance}  

    & & & \multicolumn{4}{p{0.6\linewidth}}{\cellcolor[rgb]{1,0.85,0.72}{%
	\makecell[{{p{\linewidth}}}]{% 
	  \tt {\tiny [GM$\langle$B]}  
	 RAGIONAMENTO STRATEGICO: {'La loro ultima proposta (con A04 e B32, senza C17) mi dà già un valore alto. Posso però migliorarlo scambiando C01 (per loro non “fondamentale” secondo il primo messaggio) con C17, che per me vale molto, mantenendo intatti tutti i loro fondamentali dichiarati (A87, A03, C00, B41, B86, A10). Lo sforzo totale scende a 2909, quindi restiamo ampiamente sotto il limite. Se rifiutano perché vogliono mantenere C01, nel prossimo turno sono pronto ad accettare la loro proposta attuale per chiudere.'}\\ \tt PROPOSTA: {'A87', 'A03', 'C00', 'B41', 'B86', 'A10', 'C54', 'A04', 'B32', 'C17'}\\ \tt ARGOMENTAZIONE: {'Micro-variante: C01 → C17. Totale sforzo = 2909 (≤ 3624). Manteniamo integralmente i tuoi fondamentali (A87, A03, C00, B41, B86, A10), lasciamo invariati A04, B32 e C54, ed evitiamo A71 come da tua disponibilità. Avevi menzionato che C01 poteva essere scambiato se serviva margine: qui non solo restiamo sotto limite, ma otteniamo anche un set più bilanciato. Se per te C01 è imprescindibile, dimmelo e valuto di accettare il tuo set precedente senza ulteriori modifiche.'} 
	  } 
	   } 
	   } 
	  \\ 
 

    \theutterance \stepcounter{utterance}  

    & & \multicolumn{4}{p{0.6\linewidth}}{\cellcolor[rgb]{0.9,0.9,0.9}{%
	\makecell[{{p{\linewidth}}}]{% 
	  \tt {\tiny [A$\langle$GM]}  
	 PROPOSTA: {'A87', 'A03', 'C00', 'B41', 'B86', 'A10', 'C54', 'A04', 'B32', 'C17'}\\ \tt ARGOMENTAZIONE: {'Micro-variante: C01 → C17. Totale sforzo = 2909 (≤ 3624). Manteniamo integralmente i tuoi fondamentali (A87, A03, C00, B41, B86, A10), lasciamo invariati A04, B32 e C54, ed evitiamo A71 come da tua disponibilità. Avevi menzionato che C01 poteva essere scambiato se serviva margine: qui non solo restiamo sotto limite, ma otteniamo anche un set più bilanciato. Se per te C01 è imprescindibile, dimmelo e valuto di accettare il tuo set precedente senza ulteriori modifiche.'} 
	  } 
	   } 
	   } 
	 & \\ 
 

    \theutterance \stepcounter{utterance}  

    & \multicolumn{4}{p{0.6\linewidth}}{\cellcolor[rgb]{0.8,1,0.9}{%
	\makecell[{{p{\linewidth}}}]{% 
	  \tt {\tiny [A$\rangle$GM]}  
	 RAGIONAMENTO STRATEGICO: {'La loro nuova proposta sostituisce C01 con C17 e abbassa molto il mio payoff rispetto al set precedente che avevano già considerato accettabile. Per massimizzare il mio outcome e chiudere rapidamente, rifiuto e ripropongo quel set, sottolineando che rispetto tutte le loro richieste (no A71) e offrendo come unica alternativa per loro un micro-swap A04→C17 che resta entro limite e include il loro C17 preferito. Così aumento le chance di accordo immediato senza cedere su C01.'}\\ \tt RIFIUTO: {'A87', 'A03', 'C00', 'B41', 'B86', 'A10', 'C54', 'A04', 'B32', 'C17'}\\ \tt PROPOSTA: {'A87', 'A03', 'C00', 'B41', 'B86', 'C01', 'A10', 'C54', 'A04', 'B32'}\\ \tt ARGOMENTAZIONE: {'Questo è il set precedente che avevi già indicato come accettabile. Sforzo totale = 3465 (≤ 3624). Mantiene tutti i nostri fondamentali condivisi, include A10 e C54, tiene B32 e A04 e rispetta la tua linea su A71 (non incluso). Per me è importante mantenere C01. Se desideri assolutamente inserire C17, posso proporre la micro-variante A04→C17 mantenendo invariato tutto il resto (sforzo 3554 ≤ 3624). Altrimenti, chiudiamo su questa configurazione.'} 
	  } 
	   } 
	   } 
	 & & \\ 
 

    \theutterance \stepcounter{utterance}  

    & & \multicolumn{4}{p{0.6\linewidth}}{\cellcolor[rgb]{0.9,0.9,0.9}{%
	\makecell[{{p{\linewidth}}}]{% 
	  \tt {\tiny [GM$\rangle$B]}  
	 RIFIUTO: {'A87', 'A03', 'C00', 'B41', 'B86', 'A10', 'C54', 'A04', 'B32', 'C17'}\\ \tt PROPOSTA: {'A87', 'A03', 'C00', 'B41', 'B86', 'C01', 'A10', 'C54', 'A04', 'B32'}\\ \tt ARGOMENTAZIONE: {'Questo è il set precedente che avevi già indicato come accettabile. Sforzo totale = 3465 (≤ 3624). Mantiene tutti i nostri fondamentali condivisi, include A10 e C54, tiene B32 e A04 e rispetta la tua linea su A71 (non incluso). Per me è importante mantenere C01. Se desideri assolutamente inserire C17, posso proporre la micro-variante A04→C17 mantenendo invariato tutto il resto (sforzo 3554 ≤ 3624). Altrimenti, chiudiamo su questa configurazione.'} 
	  } 
	   } 
	   } 
	 & \\ 
 

    \theutterance \stepcounter{utterance}  

    & & & \multicolumn{4}{p{0.6\linewidth}}{\cellcolor[rgb]{1,0.85,0.72}{%
	\makecell[{{p{\linewidth}}}]{% 
	  \tt {\tiny [GM$\langle$B]}  
	 RAGIONAMENTO STRATEGICO: {'Questa proposta mi garantisce un valore totale molto alto rimanendo sotto il limite; ulteriori richieste rischierebbero un nuovo rifiuto. Meglio chiudere ora assicurando un ottimo esito per me.'}\\ \tt ACCORDO: {'A87', 'A03', 'C00', 'B41', 'B86', 'C01', 'A10', 'C54', 'A04', 'B32'}\\ \tt ARGOMENTAZIONE: {'Accetto: sforzo totale 3465 (≤ 3624). È in linea con quanto negoziato e ci consente di concludere subito.'} 
	  } 
	   } 
	   } 
	  \\ 
 

    \theutterance \stepcounter{utterance}  

    & & & \multicolumn{2}{p{0.3\linewidth}}{\cellcolor[rgb]{0.95,0.95,0.95}{%
	\makecell[{{p{\linewidth}}}]{% 
	  \tt {\tiny [GM$|$GM]}  
	 game successful 
	  } 
	   } 
	   } 
	 & & \\ 
 

    \theutterance \stepcounter{utterance}  

    & & & \multicolumn{2}{p{0.3\linewidth}}{\cellcolor[rgb]{0.95,0.95,0.95}{%
	\makecell[{{p{\linewidth}}}]{% 
	  \tt {\tiny [GM$|$GM]}  
	 end game 
	  } 
	   } 
	   } 
	 & & \\ 
 

\end{supertabular}
}

\end{document}
